\documentclass[10pt]{article}
\usepackage[utf8]{inputenc}
% Language and font encodings
\usepackage[french]{babel}
\usepackage{concmath}
\usepackage[OT1]{fontenc}

% Sets page size and margins
\usepackage[letterpaper, margin=1in]{geometry}
% Useful packages
\usepackage{amsmath}
\usepackage{graphicx}
\usepackage{subcaption}
\usepackage{float}
\usepackage{booktabs}
\usepackage{svg}
\usepackage{multirow}
\usepackage{multicol}
\usepackage{enumitem}
\usepackage{listings}
\usepackage{color}
\usepackage{tikz}
\usepackage{pgfplots}
\usepackage{pgfplotstable}
\usepackage{caption}
\usepackage{wrapfig}
\usepackage{lipsum}
\usepackage{fancyhdr}
\usepackage{wrapfig}
\usepackage{indentfirst}
\usepackage{hyperref}
\usepackage{csquotes}
\usepackage{amsfonts}
\usepackage{tikz}
\usepackage{pgffor}
\usetikzlibrary{patterns}
\usetikzlibrary{intersections}
\usepackage{relsize}
\usepackage{extpfeil}
\usepackage{mathtools}
\usepackage{pgfplots}
\pgfplotsset{compat=1.18}
\usepackage{amsmath}
\usepackage{tikz}
\usetikzlibrary{arrows.meta}
%indentfirst package is used to indent the first paragraph of each section
\usepackage{indentfirst}

% Line spacing for clear readability
\usepackage{setspace}
\setstretch{1.25}

\usepackage{mathtools}
\DeclarePairedDelimiter\abs{\lvert}{\rvert}%
\DeclarePairedDelimiter\norm{\lVert}{\rVert}%
\DeclarePairedDelimiter\ceil{\lceil}{\rceil}
\DeclarePairedDelimiter\floor{\lfloor}{\rfloor}

% Custom environments (in French)
\usepackage{amsthm}            % For custom theorem-like environments
\newtheoremstyle{frenchstyle}  % Custom theorem style with no indent and new line for title
    {10pt}                     % Space above
    {10pt}                     % Space below
    {}                         % Body font (empty for normal font)
    {}                         % Indent (no indent)
    {\bfseries}                % Theorem head font
    {}                         % Punctuation after theorem head
    {\newline}                 % Space after theorem head
    {}                         % Theorem head spec (can be left empty)

\theoremstyle{frenchstyle}
\newtheorem{question}{Question}
\newtheorem{solution}{Solution}
\newtheorem{definition}{Définition}
\newtheorem{exercice}{Exercice}
% Idea developpement newtheorem
\newtheorem{idee}{Idée}

\newcommand{\R}{\mathbb{R}}
\newcommand{\rel}{\mathcal{R}}
\newcommand{\Z}{\mathbb{Z}}
\newcommand{\N}{\mathbb{N}}
\newcommand{\B}{\mathcal{B}}
% Paragraph settings

\setlength{\parskip}{1em}
\renewcommand{\thefootnote}{\roman{footnote}}
% indentation priority

\begin{document}

\begin{titlepage}
	\vspace*{9em}{\centering\Huge\usefont{T1}{qzc}{m}{it}
		DM MAT 202\par}
	\vspace*{1em}{\centering\Large\usefont{T1}{qzc}{m}{it} Devoir à la Maison Mars\par}
	\vspace*{\fill}
	% Left Bottom Author informations
	\begin{minipage}[b]{.4\textwidth}
		\raggedright
		\scriptsize
		\textit{OZEL, Furkan}\\
		\textit{20504325}\\
		\textit{Galatasaray University, 4iéme Sem.}\\
		\textit{MAT 202}\\
		\
	\end{minipage}
\end{titlepage}

\tableofcontents
\newpage

\begin{exercice}
	\addcontentsline{toc}{section}{Exercice 1}
	Trouver le domaine de définition, et dessiner les lignes de niveau et le champ de vecteurs gradient des fonctions suivantes :
	\begin{align*}
		f & : D_f \xrightarrow{\hspace*{1cm}} \R    \\
		  & (x, y) \xmapsto{\hspace*{1cm}} \ln(xy)
		\\
		\\
		h & : D_h \xrightarrow{\hspace*{1cm}} \R    \\
		  & (x, y) \xmapsto{\hspace*{1cm}} \ln(x)-y
	\end{align*}
	Dans chacun des cas, donner l’équation à la tangente à la ligne de niveau passant par $(1, 1)$, et un vecteur directeur de cette droite.
\end{exercice}

\begin{solution}
	\addcontentsline{toc}{subsection}{Solution de l'exercice 1}

	Le domaine de définition de la fonction $f$ est lié à la condition $xy > 0$ car la définition analytique du logarithme naturel est donnée par l'intégrale de $1/t$ de $1$ à $x$. Par conséquent, le domaine de définition de la fonction $f$ est donné par l'ensemble suivant :
	\[
		D_f = \{(x, y) \in \mathbb{R}^2 : xy > 0\}
	\]

	Comme on peut le comprendre intuitivement, cet ensemble se situe dans les première et troisième régions du plan analytique (les quadrants où $x$ et $y$ ont le même signe). De plus,
	\begin{enumerate}
		\item $x > 0$ et $y > 0$ : le premier quadrant
		\item $x < 0$ et $y < 0$ : le troisième quadrant
		\item $x = 0$ ou $y = 0$ : la fonction n'est pas définie (les axes sont exclus)
	\end{enumerate}

	Pour obtenir le champ de vecteurs gradient de la fonction $f$, on calcule les dérivées partielles :
	\[
		\frac{\partial f}{\partial x} = \frac{y}{xy} = \frac{1}{x} \quad \text{et} \quad \frac{\partial f}{\partial y} = \frac{x}{xy} = \frac{1}{y}
	\]
	Le gradient est donc :
	\[
		\nabla f(x,y) = \left( \frac{\partial f}{\partial x}, \frac{\partial f}{\partial y} \right) = \left( \frac{1}{x}, \frac{1}{y} \right)
	\]
	Ce champ de vecteurs est défini pour tout $(x, y) \in D_f$.

	Les lignes de niveau de la fonction $f$, fixées par $f(x,y) = k$ pour une constante $k \in \mathbb{R}$, sont données par l'équation suivante :
	\begin{align*}
		f(x,y)  & = k   \\
		\ln(xy) & = k   \\
		xy      & = e^k
	\end{align*}
	Ce sont des hyperboles. La valeur de $k$ peut être n'importe quel nombre réel. Par exemple, la courbe de niveau passant par le point $(1, 1)$ correspond à $k = \ln(1 \times 1) = \ln(1) = 0$. L'équation de cette courbe de niveau est donc $xy = e^0 = 1$. C'est une hyperbole située dans le premier et le troisième quadrant.

	Pour trouver la tangente à la courbe de niveau $xy=1$ passant par $(1, 1)$, on utilise le gradient en ce point. Le gradient est normal (perpendiculaire) à la courbe de niveau.
	\[
		\nabla f(1,1) = \left( \frac{1}{1}, \frac{1}{1} \right) = (1,1)
	\]
	Le vecteur tangent $(dx, dy) = (x-1, y-1)$ doit être orthogonal au gradient $\nabla f(1,1)$. Leur produit scalaire est donc nul :
	\[
		\nabla f(1,1) \cdot (x-1, y-1) = 0
	\]
	\[
		(1, 1) \cdot (x-1, y-1) = 1(x-1) + 1(y-1) = 0
	\]
	\[
		x - 1 + y - 1 = 0
	\]
	\[
		x + y = 2 \quad (\text{ou } y = -x + 2)
	\]
	Par conséquent, la tangente à la courbe de niveau passant par $(1, 1)$ est donnée par l'équation $y = -x + 2$. Un vecteur directeur de cette droite peut être trouvé à partir de la pente $m=-1$. Si $x$ augmente de 1, $y$ diminue de 1, donc un vecteur directeur est $(1, -1)$.

	On a calculé les données suiffantes, alors on peut dessiner la courbe de niveau et le champ de vecteurs gradient de la fonction $f$ :

	\begin{figure}[H]
		\centering
		\includesvg[width=0.45\textwidth]{ln_x_minus_y_plot.svg}
		\label{fig:transport}
	\end{figure}

	Pour la fonction $h(x,y) = \ln(x) - y$, la condition vient du logarithme naturel, qui exige que son argument soit strictement positif. Donc $x>0$. Il n'y a pas de restriction sur $y$. D'où l'ensemble de définition :
	\[
		D_h = \{(x, y) \in \mathbb{R}^2 : x > 0\} \quad \text{(le demi-plan droit)}
	\]

	Calculons le champ de vecteurs gradient de la fonction $h$ :
	\[
		\frac{\partial h}{\partial x} = \frac{1}{x}, \quad \frac{\partial h}{\partial y} = -1
	\]
	\[
		\nabla h(x,y) = \left( \frac{\partial h}{\partial x}, \frac{\partial h}{\partial y} \right) = \left( \frac{1}{x}, -1 \right)
	\]
	Ce champ de vecteurs est défini pour tout $(x, y) \in D_h$.

	Fixons $h(x,y) = k$ pour une constante $k \in \mathbb{R}$. Les lignes de niveau sont données par :
	\begin{align*}
		\ln(x) - y & = k          \\
		y          & = \ln(x) - k
	\end{align*}
	Ce sont des courbes logarithmiques décalées verticalement. La courbe de niveau passant par le point $(1, 1)$ correspond à la constante $k = h(1,1) = \ln(1) - 1 = 0 - 1 = -1$. L'équation de cette ligne de niveau est donc $y = \ln(x) - (-1) = \ln(x) + 1$.

	Calculons le gradient au point $(1,1)$ :
	\[
		\nabla h(1,1) = \left( \frac{1}{1}, -1 \right) = (1,-1)
	\]
	Ce vecteur est normal à la ligne de niveau $y = \ln(x) + 1$ au point $(1,1)$.

	Pour trouver l'équation de la tangente à cette ligne de niveau au point $(1, 1)$, on utilise le fait que le vecteur tangent $(x-1, y-1)$ est orthogonal au vecteur gradient $\nabla h(1,1)$. Leur produit scalaire doit être nul :
	\[
		\nabla h(1,1) \cdot (x-1, y-1) = 0
	\]
	\[
		(1, -1) \cdot (x-1, y-1) = 1(x-1) + (-1)(y-1) = 0
	\]
	\[
		x - 1 - y + 1 = 0
	\]
	\[
		x - y = 0 \quad \Rightarrow \quad y = x
	\]
	La tangente à la courbe de niveau passant par $(1, 1)$ est donc donnée par l'équation $y = x$. La pente de cette droite est $m=1$. Un vecteur directeur est $(1, m) = (1, 1)$.

	Finalement, on a calculé les données suffisantes, alors on peut dessiner la courbe de niveau et le champ de vecteurs gradient de la fonction $h$.
	\begin{figure}[H]
		\centering
		\includesvg[width=0.5\textwidth]{ln_xy_plot.svg}
		\label{fig:transport}
	\end{figure}
\end{solution}

\newpage

\begin{exercice}
	\addcontentsline{toc}{section}{Exercice 2}
	Soit la fonction à deux variables définie par \( f(x,y) = e^{x - y} + x + y \). On regarde la ligne de niveau 1, dénotée \( L_1 \). Vérifier que \( (0,0) \in L_1 \). Justifier qu’il existe deux intervalles ouverts \( I \) et \( J \) contenant 0, et une fonction \( \widetilde{\varphi} : J \to I \) de classe \( \mathcal{C}^1 \) telle que, pour tout \( y \in J \),
	\( (x,y) \in L_1 \) si et seulement si \( x = \widetilde{\varphi}(y) \).

	\begin{enumerate}
		\item Prouver les formules des dérivées premières et secondes de \( \widetilde{\varphi} \) :
		      \[
			      \widetilde{\varphi}' = -\frac{\frac{\partial f}{\partial y}}{\frac{\partial f}{\partial x}}
			      \quad \text{et} \quad
			      \widetilde{\varphi}'' = -\frac{1}{\frac{\partial f}{\partial x}}
			      \left( \frac{\partial^2 f}{\partial y^2}
			      + 2 \widetilde{\varphi}'(y) \frac{\partial^2 f}{\partial x \partial y}
			      + \left( \widetilde{\varphi}'(y) \right)^2 \frac{\partial^2 f}{\partial x^2} \right)
		      \]

		\item Dessiner à quoi ressemble la courbe au voisinage de \( (0,0) \).
	\end{enumerate}

\end{exercice}

\begin{solution}
	\addcontentsline{toc}{subsection}{Solution de l'exercice 2}

	Soit la fonction $f(x,y) = e^{x - y} + x + y$ et la ligne de niveau $L_1$ définie par $f(x,y) = 1$.

	Calculons $f(0,0)$ :
	\[
		f(0,0) = e^{0-0} + 0 + 0 = e^0 = 1
	\]
	Puisque $f(0,0) = 1$, le point $(0,0)$ appartient bien à la ligne de niveau $L_1 = \{(x,y) \in \mathbb{R}^2 : f(x,y) = 1\}$.

	Nous voulons montrer qu'au voisinage de $(0,0)$, l'équation $f(x,y) = 1$ définit implicitement $x$ comme une fonction de $y$, soit $x = \varphi(y)$. Nous utilisons le Théorème des Fonctions Implicites (TFI).
	\begin{itemize}
		\item La fonction $f(x,y) = e^{x-y} + x + y$ est une somme et composition de fonctions $C^\infty$ (exponentielle, polynômes) sur $\mathbb{R}^2$. Donc $f$ est de classe $C^\infty$ (et a fortiori $C^1$) sur $\mathbb{R}^2$, en particulier au voisinage de $(0,0)$.
		\item Nous avons vérifié que $f(0,0) = 1$.
		\item Calculons la dérivée partielle de $f$ par rapport à $x$ au point $(0,0)$ :
		      \[
			      \frac{\partial f}{\partial x} = \frac{\partial}{\partial x} (e^{x-y} + x + y) = e^{x-y} + 1
		      \]
		      Au point $(0,0)$, cette dérivée vaut :
		      \[
			      \frac{\partial f}{\partial x}(0,0) = e^{0-0} + 1 = 1 + 1 = 2
		      \]
		\item Puisque $\frac{\partial f}{\partial x}(0,0) = 2 \neq 0$, les conditions du TFI sont satisfaites.
	\end{itemize}
	Le TFI garantit donc l'existence d'un intervalle ouvert $J$ contenant $0$ et d'un intervalle ouvert $I$ contenant $0$, et d'une unique fonction $\varphi : J \to I$ de classe $C^1$ (et même $C^\infty$ puisque $f$ est $C^\infty$) telle que pour tout $y \in J$, on a $f(x,y) = 1 \iff x = \varphi(y)$. De plus, on sait que $\varphi(0) = 0$ (puisque $(0,0)$ est le point considéré).

	On sait que pour $y \in J$, $f(\varphi(y), y) = 1$.

	En dérivant l'équation $f(\varphi(y), y) = 1$ par rapport à $y$ en utilisant la règle de dérivation des fonctions composées :
	\[
		\frac{d}{dy} f(\varphi(y), y) = \frac{\partial f}{\partial x}(\varphi(y), y) \cdot \varphi'(y) + \frac{\partial f}{\partial y}(\varphi(y), y) \cdot 1 = 0
	\]
	Donc,
	\[
		\varphi'(y) = - \frac{\frac{\partial f}{\partial y}(\varphi(y), y)}{\frac{\partial f}{\partial x}(\varphi(y), y)}
	\]
	Calculons la dérivée partielle de $f$ par rapport à $y$ :
	\[
		\frac{\partial f}{\partial y} = \frac{\partial}{\partial y} (e^{x-y} + x + y) = -e^{x-y} + 1
	\]
	Au point $(0,0)$ (donc $y=0$ et $x=\varphi(0)=0$) :
	\[
		\frac{\partial f}{\partial y}(0,0) = -e^{0-0} + 1 = -1 + 1 = 0
	\]
	Ainsi,
	\[
		\varphi'(0) = - \frac{\frac{\partial f}{\partial y}(0,0)}{\frac{\partial f}{\partial x}(0,0)} = - \frac{0}{2} = 0
	\]

	On utilise la formule générale pour la dérivée seconde (obtenue en dérivant l'expression de $\varphi'(y)$ ou l'équation implicite une seconde fois) :
	\[
		\varphi''(y) = - \frac{1}{\frac{\partial f}{\partial x}} \left[ \frac{\partial^2 f}{\partial x^2} (\varphi')^2 + 2 \frac{\partial^2 f}{\partial x \partial y} \varphi' + \frac{\partial^2 f}{\partial y^2} \right]
	\]
	(où toutes les dérivées partielles de $f$ sont évaluées en $(\varphi(y), y)$ et $\varphi'$ est $\varphi'(y)$).
	Calculons les dérivées secondes de $f$ :
	\begin{align*} \frac{\partial^2 f}{\partial x^2} &= \frac{\partial}{\partial x} (e^{x-y} + 1) = e^{x-y} \\ \frac{\partial^2 f}{\partial y^2} &= \frac{\partial}{\partial y} (-e^{x-y} + 1) = -(-1)e^{x-y} = e^{x-y} \\ \frac{\partial^2 f}{\partial x \partial y} &= \frac{\partial}{\partial y} (e^{x-y} + 1) = -e^{x-y} \end{align*}
	Évaluons ces dérivées secondes au point $(0,0)$ :
	\[
		\frac{\partial^2 f}{\partial x^2}(0,0) = e^0 = 1, \quad
		\frac{\partial^2 f}{\partial y^2}(0,0) = e^0 = 1, \quad
		\frac{\partial^2 f}{\partial x \partial y}(0,0) = -e^0 = -1
	\]
	Maintenant, substituons les valeurs en $y=0$ (donc $\varphi(0)=0$ et $\varphi'(0)=0$) dans la formule de $\varphi''(y)$ :
	\begin{align*} \varphi''(0) &= - \frac{1}{\frac{\partial f}{\partial x}(0,0)} \left[ \frac{\partial^2 f}{\partial x^2}(0,0) (\varphi'(0))^2 + 2 \frac{\partial^2 f}{\partial x \partial y}(0,0) \varphi'(0) + \frac{\partial^2 f}{\partial y^2}(0,0) \right] \\ &= - \frac{1}{2} \left[ 1 \cdot (0)^2 + 2 \cdot (-1) \cdot 0 + 1 \right] \\ &= - \frac{1}{2} [0 + 0 + 1] \\ &= -\frac{1}{2} \end{align*}

	Utilisons le développement de Taylor de $\varphi(y)$ autour de $y=0$ à l'ordre 2 :
	\[
		\varphi(y) = \varphi(0) + \varphi'(0)(y-0) + \frac{1}{2} \varphi''(0)(y-0)^2 + o((y-0)^2)
	\]
	En substituant les valeurs trouvées :
	\[
		\varphi(y) = 0 + 0 \cdot y + \frac{1}{2} \left( -\frac{1}{2} \right) y^2 + o(y^2) = -\frac{1}{4} y^2 + o(y^2)
	\]
	Ainsi, au voisinage de $y=0$ (et donc $x=0$), la courbe de niveau $L_1$ (qui est le graphe de $x=\varphi(y)$) est approchée par la parabole d'équation $x = -\frac{1}{4} y^2$.

	Puisque $\varphi''(0) = -1/2 < 0$, la fonction $\varphi(y)$ est \textbf{concave} au voisinage de $y=0$. Le graphe $x = -\frac{1}{4} y^2$ est une parabole ouverte vers les $x$ négatifs, ce qui est cohérent avec la concavité.

	Le dessin ressemble à ceci :

	\begin{figure}[H]
		\centering
		\includesvg[width=0.5\textwidth]{parabola.svg}
	\end{figure}
\end{solution}

\newpage

\begin{exercice}
	\addcontentsline{toc}{section}{Exercice 3}
	Pour tout couple $(x, y)$ où $x \neq 0$, on défini une fonction dénotée $f_{(x,y)}$ par :
	\begin{align*}
		f_{(x,y)} & : \R \xrightarrow{\hspace*{1cm}} \R            \\
		          & t \xmapsto{\hspace*{1cm}} x^2 t^2 + (x+y)t + 1
	\end{align*}
	On pose
	\begin{align*}
		F & : D \xrightarrow{\hspace*{1cm}} \R                         \\
		  & (x,y) \xmapsto{\hspace*{1cm}} \min_{t \in \R} f_{(x,y)}(t)
	\end{align*}
	Donner le domaine de définition de $F$. Calculer $F(x, y)$ et ses dérivées partielles.
\end{exercice}

\begin{solution}
	\addcontentsline{toc}{subsection}{Solution de l'exercice 3}
	
	Pour tout couple $(x, y)$ où $x \neq 0$, on définit la fonction $f_{x,y} : \mathbb{R} \to \mathbb{R}$ par :
    \[
        f_{x,y}(t) = x^2 t^2 + (x+y)t + 1
    \]
    On pose $F(x,y) = \min_{t \in \mathbb{R}} f_{x,y}(t)$.

    La fonction $f_{x,y}(t)$ est un polynôme du second degré en $t$. Pour que cette fonction admette un minimum unique (fini) sur $\mathbb{R}$, il faut que le coefficient du terme quadratique ($t^2$) soit strictement positif.
    Le coefficient est $x^2$. La condition $x^2 > 0$ est équivalente à $x \neq 0$.
    Si $x \neq 0$, $f_{x,y}(t)$ représente une parabole dont les branches pointent vers le haut, elle admet donc un minimum unique.
    Si $x = 0$, la fonction devient $f_{0,y}(t) = yt + 1$. Si $y \neq 0$, c'est une fonction linéaire non constante qui n'admet pas de minimum sur $\mathbb{R}$. Si $y = 0$, $f_{0,0}(t) = 1$ est une fonction constante dont le minimum est 1, mais le cas $x=0$ est exclu par la définition initiale de $f_{x,y}$.

    Par conséquent, la fonction $F(x,y)$ est définie si et seulement si $x \neq 0$.
    Le domaine de définition de $F$ est donc :
    \[
        D_F = \{(x, y) \in \mathbb{R}^2 : x \neq 0\}
    \]

    Pour $x \neq 0$, le minimum de la parabole $f_{x,y}(t)$ est atteint lorsque sa dérivée par rapport à $t$ s'annule :
    \[
        \frac{df_{x,y}}{dt} = \frac{d}{dt} (x^2 t^2 + (x+y)t + 1) = 2x^2 t + (x+y)
    \]

    En posant $\frac{df_{x,y}}{dt} = 0$ :
    \[
        2x^2 t + (x+y) = 0 \implies t_{min} = -\frac{x+y}{2x^2}
    \]

    La valeur minimale est $F(x,y) = f_{x,y}(t_{min})$. Substituons $t_{min}$ dans l'expression de $f_{x,y}(t)$ :
    
	\begin{align*}
        F(x,y) & = x^2 \left(-\frac{x+y}{2x^2}\right)^2 + (x+y)\left(-\frac{x+y}{2x^2}\right) + 1 \\
               & = x^2 \frac{(x+y)^2}{4(x^2)^2} - \frac{(x+y)^2}{2x^2} + 1 \\
               & = \frac{(x+y)^2}{4x^2} - \frac{(x+y)^2}{2x^2} + 1 \\
               & = \frac{(x+y)^2 - 2(x+y)^2}{4x^2} + 1 \\
               & = -\frac{(x+y)^2}{4x^2} + 1 \\
               & = 1 - \frac{(x+y)^2}{4x^2}
    \end{align*}

    Donc, $F(x,y) = 1 - \frac{(x+y)^2}{4x^2}$ pour $x \neq 0$.

    Nous devons calculer $\frac{\partial F}{\partial x}$ et $\frac{\partial F}{\partial y}$ pour $F(x,y) = 1 - \frac{(x+y)^2}{4x^2}$.

    On peut utiliser la règle de dérivation du quotient pour le terme $\frac{(x+y)^2}{4x^2}$. Posons $u(x,y) = (x+y)^2$ et $v(x) = 4x^2$.
    $\frac{\partial u}{\partial x} = 2(x+y)$.
    $\frac{dv}{dx} = 8x$.
    \begin{align*}
        \frac{\partial}{\partial x} \left( \frac{u}{v} \right) & = \frac{\frac{\partial u}{\partial x} v - u \frac{dv}{dx}}{v^2} \\
                                                             & = \frac{[2(x+y)](4x^2) - [(x+y)^2](8x)}{(4x^2)^2} \\
                                                             & = \frac{8x^2(x+y) - 8x(x+y)^2}{16x^4} \\
                                                             & = \frac{8x(x+y) [x - (x+y)]}{16x^4} \\
                                                             & = \frac{(x+y)(-y)}{2x^3} = -\frac{y(x+y)}{2x^3}
    \end{align*}

    Alors,
    \[
        \frac{\partial F}{\partial x} = \frac{\partial}{\partial x} \left( 1 - \frac{(x+y)^2}{4x^2} \right) = 0 - \left( -\frac{y(x+y)}{2x^3} \right) = \frac{y(x+y)}{2x^3} \quad \text{(ou } \frac{xy+y^2}{2x^3} \text{)}
    \]

    On dérive par rapport à $y$, en traitant $x$ comme une constante :
    \begin{align*}
        \frac{\partial F}{\partial y} & = \frac{\partial}{\partial y} \left( 1 - \frac{(x+y)^2}{4x^2} \right) \\
                                      & = 0 - \frac{1}{4x^2} \frac{\partial}{\partial y}((x+y)^2) \\
                                      & = -\frac{1}{4x^2} [2(x+y) \cdot \frac{\partial}{\partial y}(x+y)] \\
                                      & = -\frac{1}{4x^2} [2(x+y) \cdot (0+1)] \\
                                      & = -\frac{2(x+y)}{4x^2} \\
                                      & = -\frac{x+y}{2x^2}
    \end{align*}
	
    Les dérivées partielles sont donc :
    \[
        \frac{\partial F}{\partial x} = \frac{y(x+y)}{2x^3} \quad \text{et} \quad \frac{\partial F}{\partial y} = -\frac{x+y}{2x^2}
    \]
\end{solution}

\end{document}